\documentclass{beamer}
% imprimir
% \documentclass[handout]{beamer} 
% \usepackage{pgfpages}
% \pgfpagesuselayout{4 on 1}[a4paper,landscape,border shrink=5mm]

\mode<presentation> {
  \usetheme{Warsaw}
  \setbeamercovered{transparent}
}

\usebackgroundtemplate{\includegraphics[width=\paperwidth]{format/libresoft-bg.png}}
\usepackage[spanish]{babel}
\usepackage[utf8]{inputenc}
\usepackage{graphics}
\usepackage{amssymb} % Simbolos matematicos

%% Metadatos del PDF.
\hypersetup{  
  pdftitle={Presentación del curso},
  pdfauthor={Miguel Vidal, Jose Castro},
  pdfcreator={GSyC/Libresoft},
  pdfproducer=PDFLaTeX,
  pdfsubject={Arquitectura de servidores con software libre},
}

% Inicio de la presentación
\begin{document}

% Metadatos de la presentacion
\title{Presentación del curso}
\subtitle{Arquitectura de servidores con software libre}
\institute{\{mvidal,jfcastro\}@libresoft.es} 
\author{Miguel Vidal, Jose Castro}
\date{25 de marzo de 2011}

% Cover slide
\frame{
  \maketitle
  \begin{center}
    \includegraphics[width=6cm]{format/gsyc-urjc}
  \end{center}
}

%% License slide
\begin{frame}
  \vspace{2cm}
  \begin{flushright}
    {\footnotesize \copyright{2009-2011} Miguel Vidal, Jose Castro.} \\
    \vspace{0.25cm}
    {\scriptsize Algunos derechos reservados. \\
    Esta presentación se distribuye bajo una licencia Creative Commons Reconocimiento 3.0 España}\\
    \vspace{0.10cm}
  \end{flushright}
  \begin{center}
    \href{http://creativecommons.org/licenses/by/3.0/es}{\includegraphics[width=2cm]{format/cc-by.png}} \\
    {\tiny \url{http://creativecommons.org/licenses/by/3.0/es}}
  \end{center}
\end{frame}


\section{Introducción}

\begin{frame}
\frametitle{Información general}
\textbf{Objetivo}\\
\hspace{0.3cm}Ofrecer una panorámica de las soluciones que ofrece el software libre para construir complejas arquitecturas virtualizadas.\\
\textbf{Fechas}\\
\hspace{0.3cm}Inicio: 25 de marzo 2011 -- Fin: 10 de junio 2011\\
\textbf{Horario}\\
\hspace{0.3cm}Viernes de 16:30h a 20:30h\\
\textbf{Lugar}\\
\hspace{0.3cm}Madrid On Rails\\
\textbf{Material}\\
\hspace{0.3cm}Ordenador y libro de referencia\\
\textbf{Título}\\
\hspace{0.3cm}Título propio de la URJC.\\
\hspace{0.3cm}Se obtendrá el certificado de aprovechamiento del curso.
\end{frame}

\section{Temario}
\begin{frame}
\frametitle{Clases}
\begin{enumerate}
\item \textbf{Introduccion}
\item \textbf{Seguridad}
\item \textbf{Sistemas operativos libres}
\item \textbf{Redes}
\item \textbf{Storage as a Service (StaaS)}
\item \textbf{Servicios de Internet}
\item \textbf{Virtualización I}
\item \textbf{Virtualización II}
\item \textbf{Cloud Computing}
\item \textbf{Clusters de Alta Disponibilidad}
\end{enumerate}
\end{frame}

\begin{frame}
  \frametitle{Introducción}
  \textbf{Fecha}\\
  \hspace{0.3cm}25 de marzo\\
  \textbf{Conferencia}
    \begin{itemize}
      \item \textit{Invitado:} Gregorio Robles -- Director del curso
      \item \textit{Conferencia:} Introducción al software libre
    \end{itemize}
  \textbf{Sesión}
    \begin{itemize}
      \item Metodología del curso
      \item \textit{Profesores:} Jose Castro
    \end{itemize}
    \begin{itemize}
      \item Introducción a la administración de sistemas
      \item \textit{Profesores:} Miguel Vidal
    \end{itemize}
\end{frame}

\begin{frame}
  \frametitle{Seguridad}
  \textbf{Fecha}\\
  \hspace{0.3cm}1 de abril\\
  \textbf{Conferencia}
    \begin{itemize}
      \item \textit{Invitado:} Eva Castro -- Ph.D. URJC
      \item \textit{Conferencia:} IPv6
    \end{itemize}
  \textbf{Sesión}
    \begin{itemize}
      \item Seguridad proactiva
      \item \textit{Profesores:} Pedro Coca 
    \end{itemize}
    \begin{itemize}
      \item Criptografía de clave pública
      \item \textit{Profesor:} Israel Herráiz -- Ph.D. UCM
    \end{itemize}
\end{frame}

\begin{frame}
  \frametitle{Sistemas operativos libres}
  \textbf{Fecha}\\
  \hspace{0.3cm}8 de abril\\
  \textbf{Conferencia}
    \begin{itemize}
      \item \textit{Invitado:} David Villanueva -- Ártica
      \item \textit{Conferencia:} Monitorización con Pandora
    \end{itemize}
  \textbf{Sesión}
    \begin{itemize}
      \item Sistemas operativos libres
      \item \textit{Profesores:} Miguel Vidal, Jose Castro
    \end{itemize}
\end{frame}

\begin{frame}
  \frametitle{Redes}
  \textbf{Fecha}\\
  \hspace{0.3cm}29 de abril\\
  \textbf{Conferencia}
    \begin{itemize}
      \item \textit{Invitado:} Juan José Amor -- Notario de CACert
      \item \textit{Conferencia:} Certificación con CACert
    \end{itemize}
  \textbf{Sesión}
    \begin{itemize}
      \item Redes
      \item \textit{Profesores:} Miguel Vidal, Jose Castro
    \end{itemize}
\end{frame}

\begin{frame}
  \frametitle{Storage as a Service}
  \textbf{Fecha}\\
  \hspace{0.3cm}6 de Mayo\\
  \textbf{Conferencia}
    \begin{itemize}
      \item \textit{Invitado:} Javier Turégano -- Ándago
      \item \textit{Conferencia:} La experiencia con software libre en empresas TIC
    \end{itemize}
  \textbf{Sesión}
    \begin{itemize}
      \item Storage as a Service
      \item \textit{Profesores:} Miguel Vidal, Jose Castro
    \end{itemize}
\end{frame}

\begin{frame}
  \frametitle{Servicios de Internet}
  \textbf{Fecha}\\
  \hspace{0.3cm}13 de mayo\\
  \textbf{Conferencia}
    \begin{itemize}
      \item \textit{Invitado:} Álvaro Lopez Ortega -- Líder del proyecto Cherokee
      \item \textit{Conferencia:} Servidor Web Cherokee
    \end{itemize}
  \textbf{Sesión}
    \begin{itemize}
      \item Servicios de Internet
      \item \textit{Profesores:} Miguel Vida, Jose Castro
    \end{itemize}
\end{frame}

\begin{frame}
  \frametitle{Virtualización I}
  \textbf{Fecha}\\
  \hspace{0.3cm}20 de mayo\\
  \textbf{Conferencia}
    \begin{itemize}
      \item \textit{Invitado:} Víctor Fernández -- SIA Group
      \item \textit{Conferencia:} Evolución del CPD, convergencia de almacenamiento y virtualización
    \end{itemize}
  \textbf{Sesión}
    \begin{itemize}
      \item Virtualización I
      \item \textit{Profesores:} Miguel Vidal, Jose Castro
    \end{itemize}
\end{frame}

\begin{frame}
  \frametitle{Virtualización II}
  \textbf{Fecha}\\
  \hspace{0.3cm}27 de mayo\\
  \textbf{Conferencia}
    \begin{itemize}
      \item \textit{Invitado:} Juan José Amor -- OpenSistemas \& Gnome Hispano
      \item \textit{Conferencia:} Virtualización en SPARC con LDOMs
    \end{itemize}
  \textbf{Sesión}
    \begin{itemize}
      \item Virtualización II
      \item \textit{Profesores:} Miguel Vidal, Jose Castro
    \end{itemize}
    \begin{itemize}
      \item De la virtualización al Cloud Computing
      \item \textit{Profesores:} Pedro Coca
    \end{itemize}
\end{frame}

\begin{frame}
  \frametitle{Cloud Computing}
  \textbf{Fecha}\\
  \hspace{0.3cm}3 de junio\\
  \textbf{Conferencia}
    \begin{itemize}
      \item \textit{Invitado:} Constantino Vázquez -- Desarrollador OpenNebula
      \item \textit{Conferencia:} OpenNebula
    \end{itemize}
  \textbf{Sesión}
    \begin{itemize}
      \item Tutorial de OpenNebula
      \item \textit{Profesores:} Constantino Vázquez
    \end{itemize}
\end{frame}

\begin{frame}
  \frametitle{Alta disponibilidad}
  \textbf{Fecha}\\
  \hspace{0.3cm}10 de junio\\
  \textbf{Conferencia}
    \begin{itemize}
      \item \textit{Invitado:} Flor Grinberg -- Departamento jurídico de Google
      \item \textit{Conferencia:} Aspectos legales del Cloud Computing
    \end{itemize}
  \textbf{Sesión}
    \begin{itemize}
      \item Clusters de alta disponibilidad
      \item \textit{Profesores:} Miguel Vidal, Jose Castro
    \end{itemize}
\end{frame}


\section{Metodología}
\begin{frame}
  \frametitle{Metodología}
  \begin{center}
    \textbf{Moodle} -- \url{http://moodle.libresoft.es}
  \end{center}
  \begin{itemize}
    \item Material del curso: transpas, referencias...
    \item Ejercicios
    \item Foro de dudas
    \item Noticias y avisos
  \end{itemize}
  \vspace{0.15cm}
  \begin{center}
    \textbf{Syllabus}\\
  \end{center}
  \begin{itemize}
    \item Se irá completando durante el curso
    \item Información más completa de las sesiones
    \item Ejercicios y prácticas
  \end{itemize}
\end{frame}

\begin{frame}
  \frametitle{Clases}
  \begin{enumerate}
    \item Test previo de autoevaluación {\small (no evaluable)}
    \item Conferencia de un profesional -- {\small 1 hora}
    \item Resolución de ejercicios propuestos -- {\small 30 minutos}
    \item Descanso -- {\small 15 minutos}
    \item Sesión teórica -- {\small 2 horas y 15 minutos}
    \item Ejercicios y prácticas durante la semana -- {\small no más de 1 hora}
  \end{enumerate}
\end{frame}

\begin{frame}
  \frametitle{Equipos y material}
  \begin{itemize}
    \item Puedes traer tu propio portátil
    \item WIFI para el entorno:
      \begin{itemize}
        \item ESSID:   
        \item contraseña: 
      \end{itemize}
    \item Twitter: \textit{\#casul} {\small (Curso Arquitectura Sistemas URJC Libresoft)}
    \item Canal IRC: \textit{\#casul} en irc.freenode.net
  \end{itemize}
\end{frame}

\begin{frame}
  \frametitle{Evaluación}
  \begin{itemize}
    \item Ejercicios Moodle -- 40\%
    \item Práctica final -- 50\%
    \item Entradas en el blog -- 10\% \textit{(opcional)}
    \item Participación y asistencia
  \end{itemize}
\end{frame}


\section{Contacto}
\begin{frame}
  \frametitle{Contacto}
  \textbf{Correo}\\
  \hspace{0.3cm}\{mvidal,jfcastro\}@libresoft.es\\
  \vspace{0.3cm}
  \textbf{Twitter}\\
  \hspace{0.3cm}@mvidallopez y @jfcastroluis\\
  \vspace{0.3cm}
  \textbf{IRC}\\
  \hspace{0.3cm}\textit{\#casul} en irc.freenode.net\\
  \vspace{0.3cm}
  \textbf{Moodle}\\
  \hspace{0.3cm}\url{http://moodle.libresoft.es}
\end{frame}

\begin{frame}
  \frametitle{Preguntas}
  \begin{center}
    \huge{¿Preguntas y dudas?}
  \end{center}
\end{frame}

\begin{frame}
  \frametitle{Presentación}
  \begin{center}
    \huge{Os toca presentaros}
  \end{center}
\end{frame}

% Fin del documento
\end{document}
